
\chapter{Theoretische Grundlagen und Problemanalyse}

\section{Funktionsweise von LLM-APIs}

Die Kommunikation mit LLM-APIs erfolgt typischerweise über RESTful APIs, bei denen Anfragen an zentrale Server der API Provider gesendet werden.
Die Anfragen enthalten dabei den sogenannten Prompt sowie verschiedene Parameter wie Temperatur, Maximalzahl an Tokens, usw.
Die Verbindung zur API Provider wird typerscherweise mit der Transport Layer Security (TLS) verschlüsselt und verhindert somit
Man-in-the-Middle-Angriffe und das Abhören von Dritten.
Da jedoch unverschlüsselten Daten für die Verarbeitung in den neuronalen Netzen notwendig sind,
kann der Zugriff des API-Providers auf die Daten nicht verhindert werden. (European Union Agency for Cybersecurity, 2022)
Trotz Service Level Agreements (SLAs) die die Sicherheit der Daten garantiert, entsteht ein Vertrauensproblem zwischen API-Provider und API-Konsumierer,
da die Kontrolle über die Daten bei externen Parteien liegt. (Kroll \& Zittrain, 2023) Dies stellt vorallem ein Problem dar,
wenn die API-Provider unter einer anderen regulatorischen oder wettbewerbswirksamen Umgebung als der API-Konsumierer stehen.


\section{Spezifische Risiken in der IT-Infrastruktur}
Um die Bedeutung der Daten zu verstehen, lohnt es sich diese zu kategorisieren.
Es wird unterschieden zwischen personenbezogenen Daten (Personally Identifiable Information, PII) einerseits und
geschäftskritischen Infrastrukturdaten andererseits. Personenbezogene Daten umfassen dabei alle Daten, die eine natürlich Person eindeutig identifizieren können.
Das können zum Beispiel Namen, E-Mail-Adressen, Telefonnummern, Adresse, etc. sein. Bei diesen greift die Datenschutz-Grundverordnung (DSGVO)
weswegen unachtsame Verwendung der Daten zu erheblichen Geldbußen und Reputationsschäden führen kann.
Die Kategorie geschäftskritische Infrastrukturdaten umfasst hingegen alle Daten, die für die Betriebsführung relevant sind.
Dazu zählen unter anderem IP-Adressen und Netzwerksegmentierungen, Servernamen und deren Betriebssystemversionen, Datenbank-Verbindungsstrings,
API-Schlüssel und Zugangsdaten, Konfigurationsdetails von Firewalls und Router, etc.
Im Gegensatz zu personenbezogenen Daten existieren dabei keine regulatorischen Rahmenbedingungen, die die Verwendung der Daten einschränken.
Das Risiko liegt stattdessen darin, dass die Offenlegung dieser, die Sicherheit-Infrastruktur des Unternehmens, massiv
kompromittieren kann.(Microsoft Research, 2023).